Влияние ИИ-инструментов на производительность разработки программного
обеспечения нельзя описать одним универсальным коэффициентом ускорения. В
контролируемом эксперименте с GitHub Copilot участники, завершившие ограниченную
задачу, в среднем затратили на неё на 55.8\% меньше календарного времени, чем
участники контрольной группы \cite{peng2023impact}. В рандомизированном
исследовании опытных разработчиков, работавших над задачами в знакомых зрелых
проектах с открытым исходным кодом, доступ к инструментам начала 2025 года,
напротив, увеличил скорректированное активное время реализации на 19\%
\cite{becker2025measuring}. В продолжении этого исследования отмечено вероятное
улучшение новых агентных инструментов, но одновременно выявлена более глубокая
методическая проблема: разработчики выбирали задачи с учётом доступности ИИ и
параллельно работали с несколькими агентами, поэтому длительность отдельной
задачи уже не характеризовала производительность всего процесса
\cite{metr2026uplift}. Эти результаты относятся к разным задачам, инструментам,
участникам и метрикам и потому не образуют прямого противоречия. Вместе они
показывают, что вопрос об эффективности ИИ необходимо ставить для явно заданного
способа работы, а не для модели или продукта вне организационного контекста.

Для планирования разработки этого уточнения недостаточно. Даже достоверная
оценка длительности отдельной задачи с ИИ ещё не определяет срок набора
связанных работ. Автономные интервалы нескольких агентов могут перекрываться,
но задачи конкурируют за ограниченное число потоков, следуют зависимостям,
требуют интеграции и периодически обращаются к одному разработчику за
постановкой, уточнением, проверкой или приёмкой. В результате локальное
ускорение может не изменить общий срок, если не сокращает активное ограничение,
а дополнительный агент может не создать полезного параллелизма. И наоборот,
даже режим, в котором отдельная задача выполняется дольше ручного baseline,
может сократить срок всего набора работ благодаря перекрытию автономных фаз.
Следовательно, между измеренным эффектом ИИ на задаче и календарным результатом
проекта находится самостоятельная задача организации и расписания процесса.

Настоящая работа рассматривает одну \emph{человеко-агентную ячейку}: один
разработчик выполняет заранее заданный AI-поддержанный набор программных задач
с помощью до $P$ агентных потоков и доводит каждую задачу до фиксированного
критерия приёмки. Исходным сценарием служит последовательное выполнение того же
объёма одним разработчиком без ИИ. Различаются интерактивный, делегированный и
полностью автономный режимы. В интерактивном режиме человек и инструмент
совместно ведут задачу в одном потоке; в делегированном режиме постановка и
приёмка требуют человека, а между ними возможны автономные агентные фазы; в
полностью автономном режиме обязательная человеческая фаза отсутствует.
Основной предмет статьи --- интерактивные и делегированные задачи. Полностью
автономный режим допускается как предельный случай с нулевой обязательной
человеческой работой, но обнаружение новых задач и автоматическая приёмка
результатов не моделируются.

Для делегированных задач принимается правило локального закрепления: начатая
задача удерживает назначенный агентный поток до принятия результата. Если агент
запрашивает внимание занятого разработчика, он ожидает и не получает другую
задачу, тогда как остальные агенты продолжают работу независимо. Новая задача
может начаться только в свободном потоке, а зависимая задача --- только после
приёмки всех предшественников. Такое правило описывает конкретную архитектуру
процесса, а не любые системы агентов. Переиспользование заблокированного потока,
preemption, несколько разработчиков, отдельные ограниченные ресурсы CI/API,
динамическое раскрытие графа работ и межпроектное планирование находятся за
границами постановки. Выбор полностью ручного или AI-поддержанного способа для
отдельных задач также считается внешним решением; модель планирует уже
выбранный AI-поддержанный набор.

Исследовательский вопрос формулируется следующим образом: \emph{при каких
условиях полностью заданная конфигурация работы одного разработчика с
несколькими ИИ-агентами сокращает время завершения фиксированного набора
связанных программных задач по сравнению с последовательной работой без ИИ и
какие ограничения определяют достижимое ускорение?} Конфигурация задаёт число
агентных потоков и режим каждой задачи; содержательно полное сравнение должно
также фиксировать либо явно изменять декомпозицию, граф зависимостей, фазовые
профили, критерий приёмки и функции издержек. Целевой показатель --- makespan
принятого результата. Качество и стоимость вычислений учитываются только через
их влияние на длительность проверки, исправления и завершения задачи и не
выступают самостоятельными целями оптимизации.

Основной тезис статьи состоит в том, что \textbf{достижимое ускорение является
свойством полностью заданного человеко-агентного процесса, а не произведением
``скорости ИИ'' на номинальное число агентов}. Срок может ограничиваться
суммарной загрузкой потоков, неделимой задачей, технологическим критическим
путём, последовательным человеческим ресурсом либо смешанной ресурсной цепочкой,
возникающей из порядка фаз и очереди к разработчику. Поэтому сокращение задачи,
не лежащей на текущем критическом пути, всё же может быть полезно, если оно
освобождает дефицитный поток; увеличение доступной мощности может уменьшать
срок до насыщения, но переход к конфигурации с большим настроенным $P$ способен
ухудшить его, если одновременно растут интеграционные издержки или стоимость
переключения внимания. Декомпозиция действует столь же неоднозначно: она
снимает искусственную сериализацию и открывает параллельную работу, но добавляет
границы постановки, проверки и интеграции.

Для ответа на исследовательский вопрос предложена детерминированная фазовая
модель M0--M4. Нормированная длительность $k_i^{(r)}$ относится к тройке
``задача --- инструмент --- режим'' и сравнивает полный цикл получения
принятого результата с базовой длительностью той же задачи. Это описательная
нормировка, а не причинная оценка без отдельного идентификационного дизайна. На
уровне M4 она раскладывается на автономную агентную и человеческую составляющие
$a_i^{(r)}$ и $h_i^{(r)}$, а их внутренний порядок задаётся последовательностью
фаз. Иерархия начинает с идеально делимой агрегированной работы M0 и затем
последовательно учитывает неделимость задач M1, граф зависимостей M2,
координационные издержки M3 и точное расписание фаз с общим человеческим
ресурсом и локальной блокировкой M4. Такая конструкция отделяет простые
необходимые ограничения от комбинаторных потерь конкретного расписания.

Работа вносит четыре связанных результата.

Во-первых, задана операциональная единица сравнения: не ИИ-инструмент сам по
себе, а полный сценарий ``задачи --- режимы --- ресурсы --- критерий приёмки''.
Это позволяет отделить наблюдаемую длительность задачи от прогноза, доступную
мощность от настроенной конфигурации и автономную агентную фазу от полностью
автономного режима.

Во-вторых, для уровней M0--M2 получены условия, связывающие нормированные
длительности, параллелизм и структуру работ. На M0 ускорение эквивалентно
$P>\bar{k}_w$. На M1--M2 это условие остаётся необходимым, а классическая
граница списочного планирования даёт достаточный критерий через относительную
длину самой длинной задачи или критического пути. Тем самым явно разделяются
идеальная оценка ускорения, структурная нижняя граница и достижимый срок.

В-третьих, для M4 введена структурная нижняя граница
\[
  B_4=\max\{W_4/P,L_4,\gamma(P)H\},
\]
и доказан потолок ускорения, задаваемый последовательным человеческим
ресурсом. Граница объединяет загрузку потоков,
исходный критический путь и суммарную human-work, но в общем случае не равна
оптимальному makespan $T^*$. Ресурсно-дополненный фазовый граф даёт точное
представление $T^*$ и показывает, как порядок обслуживания человека и порядок
задач на агентных слотах создают смешанные пути, отсутствующие в исходном
task-DAG. Это различает нижнюю границу, доказанный оптимум и срок конкретной
политики $T(\pi)$.

В-четвёртых, воспроизводимый вычислительный пакет с точным решателем и
фиксированной детерминированной политикой проверяет внутреннюю логику модели на
сериях EXP-00--EXP-08. Он воспроизводит иллюстративные сценарии, предъявляет
контрпримеры универсальной достижимости $B_4$, разделяет эффект доступной
мощности и растущих конфигурационных издержек, исследует гранулярность и её
взаимодействие с $P$, а также проверяет масштабную устойчивость и перенос
направлений эффектов на синтетические DAG. Эти результаты являются
вычислительной верификацией и анализом чувствительности формальной модели, а не
эмпирической оценкой конкретной команды или ИИ-инструмента.

Математические элементы этой конструкции по отдельности известны: неделимые
работы, критический путь, ограниченные ресурсы, общий обслуживающий ресурс и
дуги ресурсного порядка принадлежат теории расписаний. Заявляемый вклад состоит
не в повторном введении этих объектов, а в их последовательном отображении на
полный цикл программной задачи с task/mode-specific коэффициентом $k$,
раздельными agent/human-фазами, повторными обращениями к разработчику,
удержанием агентного потока, издержками интеграции и критерием принятого
результата. Эмпирическая калибровка такого отображения остаётся открытой
задачей.

Статья построена в соответствии с этой логикой. После постановки и компактных
обозначений сопоставляются эмпирические оценки ИИ-продуктивности, модели
масштабирования, теория расписаний и исследования человеческого внимания.
Затем вводится иерархия M0--M4 и точное ресурсное представление. Вычислительная
часть сначала фиксирует метод и правила доказательности, после чего разбирает
иллюстративный сценарий и серии EXP-00--EXP-08. Обсуждение связывает формальные
и вычислительные результаты с практическим использованием модели; отдельные
разделы фиксируют ограничения, угрозы валидности и направления развития.
Полный словарь, подробные расчёты и дополнительные результаты вынесены в
приложения.
